\documentclass[12pt]{article}
\usepackage{amsmath}
\usepackage{amsthm}
\usepackage{bm}
\usepackage[hmargin=1in,vmargin=1in]{geometry}
\usepackage[shortlabels]{enumitem}
\usepackage{multicol}
\usepackage{hyperref}
\usepackage{graphicx}
\usepackage[section]{placeins}
\usepackage{amssymb}
\usepackage{cite}
\usepackage{listings}
\usepackage{amsmath,tkz-euclide}
\usepackage{tikz}
\usetikzlibrary{arrows.meta, decorations.markings}
\usetkzobj{all}
\counterwithin{figure}{section}
\graphicspath{{/Users/edwardmcdugald/NaturalStripePatterns/presentations/figs/}}
\numberwithin{equation}{section}

\parindent 0in
\parskip 1em
\Urlmuskip=0mu plus 1mu

\title{Dissertation Roadmap}
\author{Edward McDugald}

\begin{document}

\maketitle


\section{Proposed Dissertation Objectives}
\subsection{Phase Extraction Via Autoencoder}
\par While performing phase extraction is straight forward in the traditional setup where the the input data is a pattern and the labels are the phases, the task is much harder when we do not include the phase surfaces in the loss function.
For the purpose of systems identification, this capability is essential. Thus, the first challenge we propose to tackle is to regularize an autoencoder in such a way that the latent variable is indeed the phase. There is a wealth of literature on the identification of latent space dynamics, however the interpretability of the latent space is often treated as a secondary preference. Typically, what one wants in a latent space is a significant reduction in dimensionality, so that expensive simulations can be done in the latent space, then decoded to the desired coordinate.
Our goal differs from this common approach, in that we desire strong interpretability of the latent variable, ideally as an order parameter.
\par As a first step, we propose to generate patterns whose phases are known to satisfy the self-dual RCN equation. In this way, we can add a penalty to the reconstruction loss that demands that the latent variable obey the self-dual RCN. This poses a technical challenge in learning how to access specific portions of the computational graph of the autoencoder to calculate accurate derivatives of the encoder.
Another idea that can aid this task is to pre-train the encoder in the traditional way, using a dataset of slowly varying stripe patters and exactly computed phases. The weights of this model can be used to initialize the encoder. Other interesting ideas include how to incorporate the microscopic gradient field as information to the autoencoder, as well as a wavelet based estimate of the spatial wave number. This wouldn't be too disimilar from the concept of convolutional neural nets, where one believes that a set of feature enhanced images is more informative than the original underlying image.

\subsection{Order Parameter Dynamics Identification}
\par Assuming a successful mapping to the phase from the encoder via a soft constraint (and potentially other architectural adjustments), the next step would be to combine the autoencoder with the SINDy procedure, fitting a library of candidate functions to the time derivatives of the latent variable. Again, the test case here would be patterns known to obey the RCN equations. Assuming a successful identification, we could then present the method with Swift-Hohenberg fields that have not been analytically/computationally justified from RCN. Moving further, we could consider Swift-Hohenberg fields with quadratic interactions, and solutions of Ginzburg-Landau.
Another interesting experiment would be discovery of macroscopic phase and amplitude dynamics from the Duffing oscialltor.

\subsection{Analytical and Numerical Tasks}
\par While systems identification is the prevailing theme of the proposal, there are various interesting analytical and numerical tasks that can be undertaken.
For the elliptical container, it would be desirable to understand how the point defects along the major axis of the ellipse are distributed. If this can be determined, one can define precisely the boundary conditions for convection in the ellipse, and undertake an analytical or numerical study of that scenario.
\par There is also the matter of the amplitude equations. In the analysis to date, we have stayed away from defects, allowing us to consider the first order algebraic amplitude equation. However, near defects the second order contribution to the amplitude becomes non-negligible, and this interaction has not been analytically explored.
Further, there are certain assumptions made about various terms of the macroscopic energy that allow us to ignore them from the full equation. These assumptions seem reasonable when looking at fields where the wave number is already close to its preferred value almost everywhere. However, many interesting patterns emerge when initializing SH simulations with spatial wave numbers far outside of this regime. It would be interesting to see if these previously ignored terms manifest in the dynamics when the wavenumber is far from preferred.
\par Another interesting analytical challenge is to analyze the Hilbert Transform method for slowly varying phases.
\par As a final suggestion, it would be interesting to simulate the full 3D Oberbeck-Bussinesq equations using software such as PySPH, specifically on the ellipse.

\section{Professional Development}
\par I am participating in a SIAM conference this summer of high relevance to this research. I am also continuing work with Los Alamos National Lab this summer, working on diffusion models. I have a good contact with a systems identification specialist, Youngsoo Choi, out of Lawrence Livermore. In fact, I was going to work with him this summer through the NSF-MSGI, but the program funding never came through. I am still in correspondence with him, and hope he can be of assistance in this work. I plan to continue working with the national labs, attending conferences, and exploring my options in industry, including options in finance and biotech.

\section{Timeline}
\par I hope to complete my dissertation in two years, making a total PhD timeline of 5 years. I need one more math class, which I suspect will be Sunder's topics course in stochastic PDEs. Time/energy permitting, I would be interested in taking a course in Quantum Mechanics and Statistical Mechanics as well.





\end{document}


